\section{Conclusion}\label{sec:conclusion}
% [Ky]: draft; every number is quoted from Sections 3.3 to 3.8. Adjust once the title is decided.
We compared K-means, K-means++, DBSCAN, HDBSCAN, GMM and spectral clustering on seven datasets, each
chosen to test one data property, with every method tuned by the same label-free rule. Density and graph
methods beat centroids when clusters are separated by gaps and have any shape, in few dimensions: on D1
spectral clustering recovers the moons exactly where K-means reaches ARI 0.48, and on D2 HDBSCAN follows
belts that K-means and one DBSCAN radius merge. Centroid and mixture models win when classes are compact
and touch, or the data have many dimensions: on D4 GMM reaches ARI 0.618 and K-Means++ 0.508 while
HDBSCAN puts 97\% of proteins in one cluster, and on D5 the $k$-based methods reach 0.34 to 0.40 against
0.03 or below for the density methods. Where the data have no gaps at all, as in the core of D3 and in
D8, neither family finds natural groups: the density methods return one cluster holding 92\% to 99.7\% of
the points, and K-means draws borders where its objective puts them.

The silhouette did not tell us which case we were in. On all four labelled datasets the method with the
top silhouette was not the one that matched the labels best, and on D1 it chose a DBSCAN setting that
breaks the moons into 72 fragments over one that recovers them exactly. Rerunning each result on fake
data helps where a method finds nothing real (D1 DBSCAN, D4 and D5 HDBSCAN), but not where it returns one
dominant cluster (D8) or the wrong structure (D2). The choice of method should follow from what the data
look like and what the analysis needs, and every score should be checked against fake data, against the
points it ignores, and against the share of the largest cluster.

\subsection*{Limitations}
\habegin
% [Hoang Anh]: my line for D6. Source: results/datasets/wine_results.csv (label_free_choice).
\hatag{}On D6, the methods that need $k$ were given the number of classes, which an analyst would not
know. The label-free criterion agrees for K-means (it also picks $k=3$), but BIC alone would give GMM 4
components, and D6 is small (178 points), so one wine moves a score noticeably.
\haend

% [Ky]: my line for D3 and D4. Sources: results/datasets/uber_pickups_results.csv, ecoli_results.csv.
D3 has no labels, so its clusters are judged against known geography only, and it is one 15{,}000-point
sample of one month. The D4 settings with the best ARI use the labels and overstate what an analyst could
reach without them; D4 is also small (336 proteins), so a single protein moves a rare class's share by up
to 50\%.

% [Wei Yew]: my line for D2 and D8. Source: results/datasets/{earthquakes_2023,hdb_resale}_results.csv.
D2 and D8 have no labels, so their failures are judged against plate-boundary geography and against what
the clusters contain, and cannot be scored. The density parameters are counts of points, so the clusters
would change with the size of the catalogue or the sample. The fake data shuffle each column on its own;
for the unit vectors of D2 this moves points off the sphere, so the D2 null silhouettes are a rough
baseline only.

\todo{Lam}{one line: datasets go up to 50 dimensions only (D5).}
